% =====================================================================
% Capítulo 6. Conclusiones
% Estado: redactado el 06/11/2026 (sesión R4). No introduce cifras nuevas:
% recoge las de los capítulos 4 y 5.
% =====================================================================

\chapter{Conclusiones}\label{cap:conclusiones}

Este capítulo recoge lo que el trabajo ha conseguido y lo que no. Recorre
primero los objetivos específicos con los criterios de la
tabla~\ref{tab:criterios}, resume después las conclusiones sobre la pregunta de
investigación, expone las limitaciones que acotan esas conclusiones y termina
con las líneas de trabajo futuro. La interpretación detallada de los resultados
está en la sección~\ref{sec:discusion}.

\section{Cumplimiento de los objetivos}\label{sec:concl-objetivos}

La tabla~\ref{tab:concl-objetivos} resume el resultado de cada objetivo frente
al criterio de cumplimiento fijado en el capítulo~\ref{cap:objetivos}. Como allí
se estableció, un objetivo de comparación se cumple cuando la comparación se ha
hecho correctamente, sea cual sea el modelo que resulte preferible. Todos los
objetivos se consideran cumplidos; en dos, OE8 y OE9, con una salvedad que se
explica tras la tabla.

\begin{table}[htbp]
	\centering
	\caption{Resultado de cada objetivo específico}
	\label{tab:concl-objetivos}
	\small
	\begin{tabularx}{\textwidth}{l >{\raggedright\arraybackslash}X >{\raggedright\arraybackslash}p{3.0cm}}
		\hline
		\textbf{Objetivo} & \textbf{Resultado} & \textbf{Dónde} \\
		\hline
		OE1 & Conjunto de 70\,127 horas y 60 columnas, regenerable con el código a partir de las fuentes y verificado con una huella de contenido. Sin huecos ni valores nulos; ningún valor de demanda corregido; 2927 horas de la pandemia marcadas y excluidas del entrenamiento & Secciones~\ref{sec:datos} y~\ref{sec:entorno} \\
		OE2 & Picos del periodograma en 24 y 168~h; ADF y KPSS rechazan ambos su hipótesis nula, sin contradicción, y no se diferencia la serie; relación en U con la temperatura, modelizada con grados-día. De ahí se derivan las decisiones de modelización & Sección~\ref{sec:exploratorio} \\
		OE3 & Ingenuo estacional y SARIMAX (ARMA(5,1), con términos de Fourier, exógenas, variable de nivel y retardos de 48 y 168~h) con previsiones para todo 2024. Ljung--Box rechaza el ruido blanco de los residuos: queda una dependencia diaria sin modelar & Secciones~\ref{sec:val-sarimax} y~\ref{sec:final-precision} \\
		OE4 & LightGBM y N-HiTS con la temperatura como covariable y previsiones para todo 2024; hiperparámetros finales y espacio de búsqueda documentados & Secciones~\ref{sec:esp-finales} y~\ref{sec:val-seleccion} \\
		OE5 & Protocolo con orígenes móviles, variables, horizonte y presupuesto de búsqueda comunes, 2023 para elegir y 2024 para evaluar una sola vez, y prueba de fuga con control positivo, aplicado a los cuatro modelos & Sección~\ref{sec:protocolo} \\
		OE6 & Tabla de MAE, RMSE, MAPE y MASE en 2024 y contraste de Diebold--Mariano de los seis pares. LightGBM y N-HiTS superan al SARIMAX y al ingenuo; el SARIMAX no se distingue del ingenuo & Secciones~\ref{sec:final-precision} y~\ref{sec:final-dm} \\
		OE7 & Cobertura, anchura, puntuación de intervalo y CRPS$_{19}$ de los intervalos propios y de los de ACI, con diagramas de fiabilidad, en 2023 y en 2024 & Secciones~\ref{sec:calibracion} y~\ref{sec:calibracion-final} \\
		OE8 & Precisión por hora, bloque, tipo de día y régimen, y calibración por bloque y tipo de día, con 19 contrastes. No se encuentra un régimen en el que cambie de forma distinguible el modelo preferible & Sección~\ref{sec:desglose} \\
		OE9 & Código, semillas, particiones y dependencias fijadas en un repositorio, con las tablas de resultados & Sección~\ref{sec:entorno} y anexo~\ref{anexo:repositorio} \\
		\hline
	\end{tabularx}
\end{table}

\paragraph{Salvedad de OE8.} El criterio pide las métricas desglosadas por hora
del día y por tipo de día. La precisión se presenta por hora, y la calibración
por franja de seis horas y por tipo de día, no por hora: la cobertura por hora
local se calculó y se guarda en el repositorio, pero no se ha analizado. El
puente es un solo día de 2024, y los resultados sobre él son descriptivos.
Además, el criterio incluye identificar los regímenes en que cambia el modelo
preferible; el resultado es que no cambia ninguno de forma distinguible, lo que
se informa como hallazgo.

\paragraph{Salvedad de OE9.} El criterio exige un repositorio público que permita
reproducir las tablas de la memoria. Las tablas de resultados están en el
repositorio y los datos de entrada se regeneran con el código, pero las
previsiones guardadas no se publican, de modo que la reproducción requiere
volver a entrenar los modelos, y no se ha comprobado que dos entrenamientos de
N-HiTS en equipos distintos den previsiones idénticas
(sección~\ref{sec:entorno-limites}).

\section{Conclusiones sobre la pregunta de investigación}\label{sec:concl-pregunta}

La pregunta de investigación tiene tres partes, y la
sección~\ref{sec:disc-respuesta} las responde con detalle. En síntesis, con las
condiciones de este trabajo:

\begin{enumerate}
	\item Los modelos de aprendizaje automático y profundo mejoran al SARIMAX de
	forma estadísticamente significativa en 2024. LightGBM reduce el MAE un 47\,\%
	y N-HiTS un 39\,\%, y LightGBM es apreciablemente mejor que N-HiTS. El orden
	de precisión es LightGBM, N-HiTS, SARIMAX e ingenuo, con la salvedad de que
	el SARIMAX no se distingue del ingenuo con la regla fijada de antemano. Ese es
	el resultado negativo más relevante: el modelo clásico con estacionalidad
	múltiple no mejora de forma demostrable a la referencia más simple.
	\item Lo hacen en casi todas las horas y tipos de día. LightGBM tiene el menor
	MAE en 19 de las 24 horas y N-HiTS en cinco, y ningún régimen cambia el modelo
	preferible. La mayor distancia con los modelos clásicos está en los festivos,
	donde estos fallan con un MAE de 3672 y 2813~MW frente a 796 y 655~MW, aunque
	con 10 festivos de un solo año.
	\item Ninguna familia tiene intervalos propios dentro de la banda del 88 al
	92\,\%: los de los modelos clásicos sobrecubren por poco y los cuantiles de
	LightGBM y N-HiTS infracubren claramente (76,1 y 78,7\,\%). Con los intervalos
	construidos con el mismo método, la ventana de residuos recientes o ACI, los
	cuatro quedan dentro de la banda. La calibración depende más de cómo se
	construye el intervalo que de la familia del modelo. ACI no mejora la ventana
	de residuos con el $\gamma$ elegido.
\end{enumerate}

No se concluye que una familia sea mejor en general, ni que N-HiTS sea más
robusto en épocas difíciles: esto último es una hipótesis. En cuanto a la
calibración, la construcción más sencilla, la ventana de residuos recientes sin
adaptación, da los mismos resultados que ACI con el $\gamma$ elegido, de modo
que en este problema no se justifica la complejidad añadida de ACI.

\section{Limitaciones}\label{sec:concl-limitaciones}

Las limitaciones que acotan estas conclusiones se discuten en la
sección~\ref{sec:disc-limites}. Son cinco:
\begin{itemize}
	\item \textbf{Previsión meteorológica perfecta.} Se emplea la temperatura de
	reanálisis como si fuera una previsión sin error, de modo que los resultados son
	una cota superior de lo que se obtendría en operación.
	\item \textbf{Un sistema, un año y un modelo por familia.} Las conclusiones se
	refieren al sistema peninsular, a 2024 y a los cuatro modelos estudiados, y no a
	las familias en general.
	\item \textbf{Selección con 2023.} Los hiperparámetros y el $\gamma$ de ACI se
	eligieron con 2023, y la comparación entre 2023 y 2024 mezcla el efecto del
	año, el de la elección y el del reajuste mensual, para el que no hay control.
	\item \textbf{Causas sin comprobar.} Las explicaciones del autoconsumo, del
	perfil aplanado de N-HiTS, de la infracobertura de los cuantiles, de los meses
	difíciles y del empeoramiento del SARIMAX en 2024 son hipótesis.
	\item \textbf{Especificación del SARIMAX.} Sin términos ARMA estacionales ni
	diferenciación, y con una dependencia diaria residual; no se ha explorado otra
	especificación, de modo que el resultado es el de este SARIMAX y no el de toda
	la familia.
\end{itemize}

\section{Trabajo futuro}\label{sec:concl-futuro}

Las líneas siguientes quedan fuera del alcance de este trabajo y se derivan de
sus limitaciones.
\begin{itemize}
	\item \textbf{Previsiones meteorológicas reales.} Repetir la comparación con
	las previsiones de temperatura disponibles a las 12:00 del día $D$, para medir
	cuánto se reduce la ventaja de los modelos que más dependen de la temperatura.
	Es la extensión de mayor valor práctico.
	\item \textbf{Autoconsumo y otras variables.} Incorporar una variable que
	recoja la producción solar, para comprobar si el exceso de previsión al
	mediodía desaparece.
	\item \textbf{Más sistemas, años y modelos.} Repetir el protocolo en otros años
	y sistemas, y añadir otros representantes de cada familia, como el suavizado
	exponencial con varias estacionalidades, TBATS o el \textit{Temporal Fusion
	Transformer}, que se tratan en el marco teórico pero no se implementan.
	\item \textbf{Un SARIMAX con términos estacionales.} Explorar si otra
	especificación acerca al modelo clásico a los de aprendizaje automático.
	\item \textbf{Controles que faltan.} Una evaluación sin reajuste mensual, para
	separar su efecto del del año, y una comparación de la cobertura de los
	cuantiles dentro y fuera de la muestra de entrenamiento.
	\item \textbf{Calibración.} Evaluar ACI con otros valores de $\gamma$ y otros
	métodos de calibración, y analizar la cobertura por hora local.
	\item \textbf{Lectura económica del error.} Traducir el error de previsión en
	su coste para el mercado diario, que la propuesta mencionaba y el alcance
	dejó fuera.
\end{itemize}
